\originalchapter{平动、滚动与能量}\label{ch:rolling}
\originalsection{刚体速度与无滑条件}
刚体的一般运动可分解为某点的平动加上绕该点的转动. 取质心 $C$, 任意物质点速度为
\begin{equation}
\vct v_P=\vct V_C+\vct\omega\times\vct r_{P/C}.
\end{equation}
这是瞬时关系; “瞬时转动中心”可随时间改变, 不能因此把任意时段都当作绕固定点转动. 在静止水平面上向右无滑滚动的圆轮, 取 $\omega_z<0$, 接触点速度为 $V_C+R\omega_z=0$, 所以 $V_C=R|\omega_z|$. 顶点速度为 $2V_C$, 接触点瞬时零速. 接触点仍有向上的加速度, 不是固定支点.

位于轮缘的点从顶部出发, 轮心匀速向右, 写 $\phi=Vt/R$. 其轨迹为 $x=x_0+Vt+R\sin\phi$, $y=R+R\cos\phi$, 是摆线的移相形式. 对移动表面无滑, 应令两接触物质点速度相等, 不能继续置接触点地面速度为零.

\originalsection{滚动动力学}
物体滚下倾角 $\beta$ 的固定斜面, 令向下加速度 $a$, 静摩擦 $f$ 向上, 得 $Mg\sin\beta-f=Ma$, $fR=I_C\alpha$, $a=R\alpha$. 因而
\begin{equation}
a=\frac{g\sin\beta}{1+I_C/(MR^2)},\qquad
f=\frac{I_C}{MR^2+I_C}Mg\sin\beta.
\end{equation}
实心球、实心圆柱、薄环依次有 $I_C/(MR^2)=2/5,1/2,1$, 因而同坡上的加速度依次为 $5/7,2/3,1/2$ 倍 $g\sin\beta$. 摩擦必须满足 $f\le\mu_sMg\cos\beta$, 否则该无滑解失效.

质心平动和绕质心转动的动能为 $K=MV_C^2/2+I_C\omega^2/2$. 对静止刚性斜面, 无滑接触点速度为零, 静摩擦瞬时总功率为零. 它可减少平动能增长并增加转动能, 两部分功率抵消. 因而高度差 $h$ 给 $Mgh=(M+I_C/R^2)v^2/2$. 移动地面上的静摩擦可以做功, 不可笼统说静摩擦永不做功.

\originalsection{滑动到滚动}
物体若初始 $v_0$ 与 $R|\omega_0|$ 不匹配, 接触点有相对滑动. 取向右为正、逆时针角速度为正, 接触速度 $u=v+R\omega$. 若 $u>0$, 动摩擦向左, $\dot v=-\mu_kg$, $\dot\omega=-\mu_kMgR/I_C$. 两式使 $u$ 减少. 在 $u=0$ 时转入无滑关系, 不能让摩擦继续保持原大小而使接触点反向滑动.

设 $I_C=\kappa MR^2$, 则达到无滑的时间为 $t_*=(v_0+R\omega_0)/[\mu_kg(1+1/\kappa)]$ (此式针对 $u_0>0$). 末速 $v_*=(v_0-\kappa R\omega_0)/(1+\kappa)$. 转入滚动过程中耗散功为接触摩擦乘以相对滑移积分, 不能只用质心位移计算总摩擦功.
\originalsection{刚体分解的几何理由}
在很短时间内, 刚体点集的运动保持任意两点距离. 去除一个参考点的共同平移后, 剩余变换保持原点及长度、夹角, 在不反射的连续运动中是旋转. 微小转角 $\dd\vct\theta$ 使相对位置变化 $\dd\vct r=\dd\vct\theta\times\vct r$, 除以时间得到 $\vct\omega\times\vct r$. 因而第\ref{ch:rolling}章速度式不是额外假定. 有限刚体位移也可表示为平移加一次绕某轴的转动 (Chasles 定理), 但有限转角的相加一般不可交换; 瞬时角速度向量可以相加, 不能把有限旋转随意当作同样的向量运算.
